\begin{table}[htbp]
\centering
\begin{threeparttable}
\caption{Conditional and counterfactual realised-asymmetry comparisons, S\&P 500, 2017.}
\label{tab:rv-ra-conditional}
\scriptsize
\setlength{\tabcolsep}{3pt}
\begin{tabular}{@{}lrrr@{}}
\toprule
Estimand & Estimate & \multicolumn{1}{c}{95\% interval} & Origins \\
\midrule
@@ROW_0@@
@@ROW_1@@
@@ROW_2@@
@@ROW_3@@
@@ROW_4@@
@@ROW_5@@
@@ROW_6@@
@@ROW_7@@
@@ROW_8@@
\bottomrule
\end{tabular}
\begin{tablenotes}[flushleft]
\footnotesize
\item Notes: Intervals are paired circular moving-block percentile intervals based on 2,000 replications and block length 10. Negative estimates favour the first forecast or model set named in the contrast. The full panel contains 241 origins. The high-asymmetry indicator identifies 49 origins and varies within 11 monthly refits. The high-log-RV contrast is not estimable because no origin exceeds the refit-specific high-state threshold in the calm 2017 sample. Counterfactual asymmetry comparisons use the five- and ten-day horizons.
\end{tablenotes}
\end{threeparttable}
\end{table}
